\maketitle

% --- 前言 ---
\noindent\textbf{写在前面：} \\
这份文档演示了 \texttt{elegantexam} 模板的所有功能。本模板最大的特点是\textbf{审美在线}（内置设计师配色）和\textbf{操作简单}（傻瓜式开关）。请大家在整理真题时，严格按照以下规范录入。

% =======================================================
\section{色彩与主题切换}
% =======================================================

我们拒绝“辣眼睛”的自定义配色，而是锁定了 6 款符合学术审美的颜色。

\begin{enumerate}
    \item 现在的解析框、水印、标题是什么颜色？ \kh[A]
    \xx{跟随主题色}{总是蓝色}{总是黑色}{随机颜色}
    
    \begin{solution}
        \textbf{解析：} 当前设置的是 \texttt{blue}。你可以尝试在 \texttt{main.tex} 第 15 行将 \texttt{blue} 改为 \texttt{orange} 或 \texttt{red}，然后重新编译。你会发现：
        \begin{itemize}
            \item 标题颜色变了；
            \item 下划线颜色变了；
            \item 解析框的背景色自动变成了对应颜色的淡色版本；
            \item 全屏水印的颜色也跟着变了。
        \end{itemize}
    \end{solution}
\end{enumerate}

% =======================================================
\section{选择题排版规范 (三种模式)}
% =======================================================

为了版面整洁，我们提供了三种排版命令，请根据**选项文字长度**灵活选择。

\begin{enumerate}
    \item \textbf{【短选项 - 四栏】} 选项是数字或短语时，用 \texttt{\textbackslash xx}：\kh[A]
    \xx{短路}{开路}{电压源}{电流源}

    \item \textbf{【中等选项 - 双栏】} 选项约半行长时，用 \texttt{\textbackslash xxtwo}：\kh[B]
    \xxtwo{这是稍微长一点的选项A}{这是稍微长一点的选项B}
          {这是稍微长一点的选项C}{这是稍微长一点的选项D}

    \item \textbf{【长选项 - 单栏】} 选项是整句或包含长公式时，\textbf{必须}用 \texttt{\textbackslash xxone}：\kh[C]
    \xxone{戴维宁等效电阻等于端口开路电压除以短路电流}
          {这是一个非常长的选项，如果强行挤在两栏里会非常难看}
          {对外电路而言，线性有源二端网络可等效为电压源与电阻串联}
          {选项里有公式 $\oint \vec{E} \cdot d\vec{l} = 0$ 时也要用单栏}
\end{enumerate}

% =======================================================
\section{智能填空题规范}
% =======================================================

填空题命令 \texttt{\textbackslash tk\{答案\}} 会自动处理下划线长度。

\begin{enumerate}
    \item \textbf{普通填空}：若电压有效值为 220V，则最大值为 \tk{311} V。
    \item \textbf{长填空}：基尔霍夫电流定律体现了 \tk{电荷守恒定律}。
    \item \textbf{指定宽度}：如果想给刷题版留更多位置，可以指定宽度：\tk[4cm]{阻抗变换}。
\end{enumerate}

% =======================================================
\section{大题与解析录入}
% =======================================================

\textbf{1. (10分) 请简述如何导出“学生刷题版”PDF？}

\begin{solution}
    \textbf{操作步骤：}
    \begin{enumerate}
        \item 打开 \texttt{main.tex} 文件；
        \item 找到“功能开关控制台”区域；
        \item 将 \texttt{\textbackslash showanswerstrue} 修改为 \texttt{\textbackslash showanswersfalse}；
        \item 点击 Recompile，此时答案消失、解析框隐藏；
        \item 下载 PDF，重命名为“刷题版.pdf”。
    \end{enumerate}
\end{solution}

\vspace{2cm} % 手动留白示例

\textbf{2. (10分) 如何在解析中插入公式和图片？}

\begin{solution}
    直接写 LaTeX 代码即可，解析框会自动包裹它们。
    $$ P = UI \cos \varphi $$
    图片请使用 \texttt{center} 环境。
\end{solution}